% ============================================================================
% Section 8.1 — Populations and Samples (OK/TX: experimental design emphasis)
% CCSS 7.SP.A.1
% ============================================================================
\section{Populations and Samples}

\begin{stepGoal}
\begin{itemize}
    \item Tell the difference between a population and a sample
    \item Design a simple experiment with a data-collection plan
\end{itemize}
\end{stepGoal}

\begin{stepVocabBox}
    \vocabItem{Population}{The \textbf{entire} group you want to learn about.}
    \vocabItem{Sample}{A \textbf{smaller} group chosen from the population.}
    \vocabItem{Random Sample}{Every member has an \textbf{equal chance} of being chosen.}
    \vocabItem{Bias}{When a sample does \textbf{not} fairly represent the population.}
\end{stepVocabBox}

\begin{stepsCard}[How to Design a Fair Experiment and Choose a Sample]
    \stepItem{Ask a clear \textbf{question} — what do you want to find out?}
    \stepItem{Identify the \textbf{population} — the whole group you are studying.}
    \stepItem{Choose a \textbf{sampling method} — use random selection so every member has an equal chance.}
    \stepItem{Collect the \textbf{data} — survey or measure your sample.}
    \stepItem{\textbf{Analyze and conclude} — what does the data tell you about the population?}
\end{stepsCard}

% --- Worked Example 1 ---
\begin{stepExample}{A school has $600$ students. The principal wants to know their favorite sport. Design a fair experiment.}
    \stepShow{1} Question: What sport is most popular among students?

    \stepShow{2} Population: all $600$ students.

    \stepShow{3} Method: Use a random number generator to select $60$ student IDs.

    \stepShow{4} Collect: Give the $60$ selected students a one-question survey.

    \stepShow{5} Analyze: Tally the results — the sport with the most votes is likely the favorite school-wide.
    \stepResult{A random sample of $60$ from $600$ avoids bias and produces a fair result.}
\end{stepExample}

% --- Worked Example 2 ---
\begin{stepExample}{A teacher surveys the first $10$ students who arrive each morning about homework time. Is this a good sample?}
    \stepShow{1} Question: How much time do students spend on homework?

    \stepShow{2} Population: all students in the class.

    \stepShow{3} Method: first $10$ to arrive — this is \textbf{not} random (early arrivals may have different habits).

    \stepShow{4} Data collected may not represent everyone.

    \stepShow{5} Conclusion: The results are likely biased. A better method would be drawing $10$ names from a hat.
    \stepResult{The sample is biased — picking only early arrivals does not give every student an equal chance.}
\end{stepExample}

\watchOut{A large sample that is not random is still biased. Randomness matters more than size!}

% ============================================================================
% PRACTICE
% ============================================================================
\newpage
\begin{stepPractice}{Populations and Samples Practice}
    \resetProblems

    \stepPracticeHeader[step1Color]{Ask the Question and Identify the Population}
    \prob You want to know the favorite lunch of $400$ students. What is the population? \answerBlank[3cm]
    \answer{All $400$ students}

    \stepPracticeHeader[step2Color]{Evaluate the Sampling Method}
    \prob A factory tests $50$ light bulbs from a batch of $5{,}000$ using a random number table. Is this a good method? \answerBlank[3cm]
    \answer{Yes — random selection gives every bulb an equal chance}

    \prob A store surveys only customers who shop on Saturday mornings. Is this random? \answerBlank[3cm]
    \answer{No — it excludes people who shop at other times, creating bias}

    \stepPracticeHeader[step3Color]{Design an Experiment}
    \wordProblem{Design an experiment to find out which school supply (pencil, pen, or marker) is most popular among $300$ students. List all five steps.}{~}
    \answer{1) Question: Which supply is most popular? 2) Population: all $300$ students. 3) Method: randomly select $50$ students. 4) Collect: ask each their preference. 5) Analyze: tally results and find the most common choice.}

    \wordProblem{A company surveys only its own employees about a new product. Explain why this is a problem and suggest a better method.}{~}
    \answer{Employees may be biased by loyalty. The population is all potential customers. A better method is to randomly survey shoppers in several stores.}
\end{stepPractice}
